\section{从课程到工程：Agent 系统实施清单}

前面各案例看似分散，但可以收敛成一套实施路径。下面给出一份可直接用于项目 kickoff 的 checklist，帮助团队把 ``概念正确'' 落到 ``系统可交付''。

\subsection{阶段化建设路线}
\begin{table}[H]
\centering
\begin{tabular}{p{2.5cm}p{4.8cm}p{4.3cm}}
\toprule
\textbf{阶段} & \textbf{必做事项} & \textbf{验收信号} \\
\midrule
定义阶段 & 明确 observe/decide/act 闭环与环境边界 & 关键状态与动作可枚举 \\
基线阶段 & 建立可运行的分层管道（P-C-A） & 能定位失败在何层出现 \\
训练阶段 & 设计场景分布与反馈体系（human + env） & 策略改进曲线稳定上升 \\
评测阶段 & 引入多体互动与异常扰动评测 & 非 IID 情况下性能不过度塌缩 \\
治理阶段 & 纳入安全/公平/合规约束与审计日志 & 违规率可量化并可追踪回放 \\
\bottomrule
\end{tabular}
\caption{Autonomous Agent 项目实施五阶段}
\end{table}

\subsection{常见失败模式与修复策略}
\begin{knowledgebox}{失败模式 A：只做离线 benchmark}
症状：离线分数高，但线上多步任务失败率高。
修复：补充交互式评测，强制覆盖恢复路径与协同路径。
\end{knowledgebox}

\begin{knowledgebox}{失败模式 B：忽视队友/对手分布漂移}
症状：固定环境里表现稳定，换协作方后性能骤降。
修复：把 ad hoc teamwork 思想引入训练，加入未知队友策略扰动。
\end{knowledgebox}

\begin{knowledgebox}{失败模式 C：安全策略完全后置}
症状：模型本体仍产生危险决策，外层拦截频繁触发。
修复：把安全与规则奖励前置到训练期，降低策略与守卫冲突。
\end{knowledgebox}

\begin{importantbox}{一句话原则}
\textit{先保证系统可诊断，再追求策略最优；先保证行为可治理，再追求局部指标极致。}
\end{importantbox}

\subsection{证据索引：课程片段与结论对照}
\begin{table}[H]
\centering
\begin{tabular}{p{2.6cm}p{4.9cm}p{4.1cm}}
\toprule
\textbf{时间戳} & \textbf{课程片段} & \textbf{提炼结论} \\
\midrule
00:08--00:10 & RoboCup 多机器人自主对抗 & 多体协作必须在实时闭环下验证 \\
00:12--00:13 & reservation-based intersection & 协议化协调优于局部最优 \\
00:21--00:23 & TAMER 与 implicit feedback & 人类反馈可作为可训练信号 \\
00:31--00:33 & Slack 任务设定与难点 & 低保真结构学习可降真实成本 \\
00:53--00:56 & GT Sophy 策略涌现与 etiquette & 数据配比与规则设计决定上限 \\
\bottomrule
\end{tabular}
\caption{关键视频证据与笔记结论映射}
\end{table}

\subsection{本章小结}
把 Stone 的课程真正落地，需要把算法、系统、评测、治理四条线并行推进。任意一条缺失，都会导致 ``可演示但不可部署''。

